\documentclass{umk-matematika}
\begin{document}

\begin{center}
{\Large\bfseries Практическая работа 15}
\end{center}
\vspace{0.5cm}

\textbf{Задача 1.} Сформулируйте признак перпендикулярности двух плоскостей.\par\vspace{10pt}
\textbf{Задача 2.} Сформулируйте следствие из признака перпендикулярности плоскостей.\par\vspace{10pt}
\textbf{Задача 3.} В кубе $ABCDA\_1B\_1C\_1D\_1$ перпендикулярны ли плоскости ABB$_1$A$_1$ и ABCD?\par\vspace{10pt}
\textbf{Задача 4.} Приведите пример перпендикулярных плоскостей в строительстве.\par\vspace{10pt}
\textbf{Задача 5.} Плоскость $\alpha$ проходит через катет BC прямоугольного треугольника ABC ($\angle C = 90^\circ$). Катет AC образует с плоскостью $\alpha$ угол 30$^\circ$. Найдите угол между плоскостью $\alpha$ и плоскостью треугольника.\par\vspace{10pt}
\textbf{Задача 6.} Через сторону AD прямоугольника ABCD проведена плоскость $\alpha$ так, что проекция BC на $\alpha$ равна 5 см. Найдите расстояние от BC до $\alpha$, если AB = 12 см.\par\vspace{10pt}
\textbf{Задача 7.} Даны две перпендикулярные плоскости $\alpha$ и $\beta$. Прямая a лежит в плоскости $\alpha$ и параллельна линии пересечения плоскостей. Перпендикулярна ли прямая a плоскости $\beta$?\par\vspace{10pt}
\textbf{Задача 8.} В прямоугольном параллелепипеде $ABCDA\_1B\_1C\_1D\_1$ AB = 6 см, AD = 8 см, AA$_1$ = 10 см. Найдите угол между плоскостями AB$_1$C$_1$ и ABC.\par\vspace{10pt}
\textbf{Задача 9.} Плоскости $\alpha$ и $\beta$ перпендикулярны. Точка A удалена от $\alpha$ на 6 см, от $\beta$ на 8 см, от линии их пересечения на x см. Найдите x.\par\vspace{10pt}
\textbf{Задача 10.} Две стены здания перпендикулярны. На одной стене на высоте 3 м от пола закреплён кронштейн, на другой стене на высоте 4 м — другой кронштейн. Расстояние между проекциями кронштейнов на пол (вдоль линии пересечения стен) равно 12 м. Найдите расстояние между кронштейнами.\par\vspace{10pt}
\newpage
\section*{Ответы}

\begin{enumerate}
  \item Ответ: признак перпендикулярности плоскостей.
  \item Ответ: следствие из признака.
  \item Ответ: да.
  \item Ответ: стена и пол.
  \item Ответ: 60$^\circ$.
  \item Ответ: $\sqrt{119}$ см.
  \item Ответ: нет, прямая a параллельна плоскости $\beta$.
  \item Ответ: 45$^\circ$.
  \item Ответ: $x = 10$ см.
  \item Ответ: 13 м.
\end{enumerate}
\end{document}
